\begin{defi}
    Data una macchina di Turing $\mu$ di alfabeto $A$ e insieme degli stati $Q$
    \[
    \mu: 
    Q\times  A' 
    \to Q' 
    \times 
    A' 
    \times 
    \{-1,+1\}
    \]
    con $A'=A \cup \{ \square\}$, $Q'=Q \cup \{ q_F^0, q_F^1\}$:
    \begin{enumerate}
        \item  sia
\[
C = (c_0,\dots,c_n)
\]
una computazione di $\mu$, il \textit{costo in tempo di $C$} è dato dalla lunghezza di $C$ diminuita di $1$;

\item il \textit{costo in spazio} è dato dalla cardinalità dell’insieme di interi $p$ tali che almeno per una delle configurazioni $c_i = (f_i,p_i,q_i)$ si ha $p= p_i$:
\[
|\{
p \, | \, (f,p,q) \in C
\}|;
\]

\item  si dirà che $\mu$ si arresta in tempo $t$ (rispettivamente spazio $s$) sull’input $w$, se $C$ è la computazione di
$\mu$ a partire dalla configurazione iniziale associata a $w$, $C$ raggiunge una configurazione terminale e il
costo in tempo di $C$ è $t$ ( rispettivamente il costo in spazio è $s$); in questo caso si userà la notazione
$T_\mu(w):A^* \to \mathbb N$ chiamata
\textit{funzione di tempo di arresto} 
(rispettivamente $S_\mu(w)$
$S_\mu(w):A^* \to \mathbb N$ chiamata
\textit{funzione di spazio di arresto} 
) per indicare il tempo di arresto $t$ (rispettivamente lo spazio di arresto
$s$).
    \end{enumerate}
\end{defi}
Di solito ci si interessa al comportamento nel caso peggiore:
\[
T_{max}(n)
:=
\operatorname{max}_{w\in A^*, \,|w|\leq n}T_\mu(w)
\]
Si noti che la definizione di costo in tempo di una computazione associata ad una parola da una macchina
di Turing si basa sulla lunghezza della computazione e pertanto si adatta a qualsiasi algoritmo formale (non
necessariamente per una macchina di Turing). Viceversa, per quanto riguarda la nozione di costo in spazio
la definizione è basata sulla specifica configurazione e pertanto non necessariamente è trasportabile nel caso
di altri algoritmi formali.
\begin{defi}
    Sia $\mu$ una macchina di Turing sull’alfabeto $A$ e con insieme di stati $Q$, e siano $T$, ed $S$ due
funzioni
$$T : \mathbb N \to \mathbb N,$$ $$S : \mathbb N \to \mathbb N.$$
Si dice che $\mu$ si \textit{arresta in tempo T} (rispettivamente \textit{in spazio S}) se per ogni $w\in A^*$ di lunghezza al più
uguale ad $n$ la macchina si arresta sull’input $w$ in tempo (risp. spazio) al più uguale a $T(n)$ (risp. $S(n)$).
\end{defi}
Infatti, quello che vorremmo è che
\[
T_\mu(w)\leq T(n), \quad \forall w \text{ tale che } {|w|\leq n}
\]
\begin{defi}

Siano $T$ e $S$ funzioni come nella definizione sopra.
\begin{enumerate}
    \item La macchina $\mu$ di alfabeto $A$ \emph{decide l'insieme $X$ in tempo $T$} (risp. \emph{spazio $S$}) se $\mu$ decide $X$ e se esiste un intero $N$ tale che per ogni intero $n \geq N$ e per ogni parola $w \in A^*$ di lunghezza al più uguale a $n$ la macchina $\mu$ si ferma a partire da $w$ in tempo al più uguale a $T(n)$ (risp. spazio al più uguale ad $S(n)$).

\item  La \textit{classe di complessità} $\mathrm{DTIME}_1^A(T)$ (risp. $\mathrm{DSPACE}_1^A(S)$) è l'insieme che contiene gli insiemi $X \subset A^*$ per i quali esiste almeno una macchina di Turing di alfabeto $A$ che decide $X$ in tempo $T$ (risp. in spazio $S$).
\end{enumerate}
\end{defi}
Con abuso di notazione, spesso scriveremo $T(n)$ invece che $T$.
\begin{rema}
    L'indice $1$ in basso utilizzato in questa notazione fa riferimento al fatto che nella macchina viene utilizzato un solo nastro.
 Con questa notazione si può notare che l'insieme $X$ degli interi divisibili per $9$ (come pure tutti gli insiemi decidibili per automa finito) appartengono alle classi $\mathrm{DTIME}_1^{\{0,1,\ldots,9\}}(n)$ e $\mathrm{DSPACE}_1^{\{0,1,\ldots,9\}}(n)$.
\end{rema}
È anche possibile interessarsi al caso medio, allora si prenderebbe in considerazione non la funzione $T_{max}$
ma una media su tutti i possibili input
\[
T_{avg} (n) :=
\frac{\sum_{w\in A^N}T_\mu(w)}{|A^N|}
\]
dove $A^N=\{ w\in A^* \,  |  \, |w| \leq N\}$.
\begin{defi}
 Una funzione
\[
T : \mathbb N \to \mathbb N
\]
viene detta \textit{funzione di complessità} se è una funzione non decrescente e se $T(n) \geq  n$ definitivamente.
\end{defi}
\subsection{Look-up Table}
Il lemma seguente mostra che il valore puntuale delle funzioni di complessità non corrisponde alla reale
difficoltà del problema, in quanto, data una macchina di Turing che decide un problema X con un certo
tempo di arresto, è sempre possibile trovarne un’altra che decide lo stesso problema con complessità lineare
per tutte le parole di input fino ad una lunghezza fissata, dunque la nozione di complessità ha senso come
misura della difficoltà di un problema solo quando si considera il comportamento asintotico dell’algoritmo.
\begin{lemm}[Look-up Table]
Dato un insieme $X \subset A^*$, supponiamo che esista una macchina di Turing $\mu$ che lo decide in tempo $T(n)$ e spazio $S(n)$.
Per un intero qualsiasi $N$ siano
\[
T'(n) =
\begin{cases}
n & n < N \\[4pt]
T(n) + 2N & \text{altrimenti}
\end{cases}
\]
e sia
\[
S'(n) =
\begin{cases}
n & n < N \\[4pt]
S(n) & \text{altrimenti}
\end{cases}
\]
allora esiste una macchina di Turing $\mu'$ che decide $X$ in tempo $T'(n)$ e spazio $S'(n)$.

\begin{proof}
Si può modificare la macchina $\mu$ per includere nell'insieme degli stati direttamente la risposta che corrisponde ad una data parola di input (in un numero finito di casi).

Ed è quello che faremo nel caso delle parole di lunghezza inferiore uguale ad $N$, in questo modo la nuova macchina si fermerà non appena sarà stata letta la parola di input. In questo caso la macchina dovrà leggere il nastro e quindi si arresterà in tempo al più $N$. In caso fossimo in presenza di una parola più lunga bisognerebbe riprendere la computazione della macchina $M$, reimpostando la macchina nello stato iniziale di $M$ e nella posizione iniziale.

Sia
\[
Q' = (Q\setminus \{ q_0\}) \cup \{q_w \mid w \in A^N\} \cup \{q_R, q_0^\mu\}
\]
dove lo stato $q_R$ sarà utilizzato come lo stato utilizzato per il ritorno alla casella iniziale nel caso in cui l'input si trovi più lungo di $N$ simboli.

Dunque la macchina di Turing $\mu'$ è definita a partire dalla macchina $\mu$:
\[
\mu'(q,a) =
\begin{cases}
\mu(a,q) & \text{se } q \in (Q\setminus\{q_0 \} )\cup \{q_0^\mu \}, \\[4pt]
(q_{wa}, a, +1) & \text{se } q = q_w \text{ per qualche } w \in A(m) \text{ con } m \le N \text{ e } |wa| \le N, \\[4pt]
(q_R, a, +1) & \text{se } q = q_w \text{ per qualche } w \in A(m) \text{ con } a \ne \square \text{ e } |w| = N, \\[4pt]
(q_R, a, -1) & \text{se } q = q_R \text{ e } a \ne \square, \\[4pt]
(q_0^\mu, a, +1) & \text{se } q = q_R \text{ e } a = \square, \\[4pt]
(q_F^1, a, +1) & \text{se } q = q_w \text{ per qualche } w \in A(m) \text{ con } m \le N,\ a = \square \text{ e } w \in X, \\[4pt]
(q_F^0, a, +1) & \text{se } q = q_w \text{ per qualche } w \in A(m) \text{ con } m \le N,\ a = \square \text{ e } w \notin X;
\end{cases}
\]
Lo stato iniziale $q_0'$ di $\mu'$ coincide con lo stato associato alla parola vuota $q_\epsilon$ mentre lo stato $q_0$ di $\mu$ viene pensato come lo stato aggiuntivo $q_0^\mu$.
Per ogni parola di lunghezza inferiore o uguale ad $N$ la computazione viene effettuata in modo simile alla computazione per automa finito, per le altre parole la macchina $\mu'$ compie una lettura ed un ritorno alla posizione iniziale che non altera l'input e dunque quando la macchina torna alla posizione iniziale e passa nello stato iniziale (di $\mu$) la computazione si svolge esattamente come per $\mu$.
\end{proof}
\end{lemm}
\textcolor{red}{Aggiugni grafico su dimostrazione.}
\begin{rema}
    Il nome "Look-up Table" nasce per il seguente motivo.
    La macchina che si costruisce 
    nella dimostrazione
    è la macchina di partenza ma che classifica
    già la decidibilità delle parole
    di lunghezza $N$.
    Poprio come se la macchina,
    alla prima lettura della parola, dovesse controllare in una "tavola"
    se la parola c'è (allora è accettata) o meno (allora viene scartata).
    Se la parola invece è più lunga di $N$,
    una volta che vengono letti 
    i primi
    $N$ caratteri,
    la macchina capisce che ci sono altri caratteri e allora,
    facendo altri $N$
    passi, torna all'inizio della parola e comincia ad analizzare la parola usando la strategia
    della macchina di Turing di partenza.
\end{rema}

\begin{rema} Il significato di questo lemma è che ha senso analizzare la complessità (dell'algoritmo di decisione associato), soltanto relativamente ad un insieme $X$ infinito, poiché nel caso di un $X$ finito si può sempre decidere il problema in un numero di passi finito, pari alla lunghezza della sua codifica sull'alfabeto (associando alla parola uno stato ``ad hoc'' e facendo una transizione accettante in base alla parola). Inoltre, poiché è possibile decidere il complementare di un insieme decidibile (invertendo nella tabella lo stato accettante $q_F^1$ e quello non accettante $q_F^0$), in pratica la questione della decidibilità si pone solo quando $X$ è infinito e co-infinito, mentre è banale quando l'insieme è finito oppure co-finito.
\end{rema}

\begin{rema} Il modello di calcolo ``macchina di Turing'' è più potente del modello ``automa a stati finiti'', in quanto per ogni automa $(T, q_0, F)$ resta definita una macchina di Turing $\mu$ sullo stesso alfabeto e sullo stesso insieme di stati per cui la computazione dell'algoritmo associato alla macchina $\mu$ coincide con l'algoritmo associato all'automa, più precisamente è sufficiente prendere
\[
\mu(q,a) :=
\begin{cases}
(T(q,a), a, +1) & a \ne \square \\[4pt]
(q_F^1, \square, +1) & \text{se } q \in F \\[4pt]
(q_F^0, \square, +1) & \text{se } q \notin F
\end{cases}
\]Ogni macchina di Turing associata ad un automa finito si arresta in tempo e spazio $T(n) = n$.    
\end{rema}
\begin{prop}
Sia $T : \mathbb{N} \to \mathbb{N}$ una funzione di complessità le inclusioni seguenti sono sempre vere.

\[
\mathrm{DTIME}_1^A(T) \subset \mathrm{DSPACE}_1^A(T) \subset \bigcup_{\epsilon > 0} \mathrm{DTIME}_1^A(2^{\epsilon T}).
\]

\begin{proof}

    Fissato l'alfabeto $A = \{0,1\}$, sia $X \in \mathrm{DTIME}_1^A(T)$. Poiché ad ogni passo di computazione è accessibile una sola posizione allora la massima posizione raggiungibile in $T(n)$ passi di computazione non può essere maggiore di $T(n)$.

Sia ora $\mu$ una macchina di Turing di alfabeto $A$ ed insieme degli stati $Q$, e sia $k = |A|$ e $d = |Q|$.

Per ogni intero $s$ esistono $(2s+1)k^{2s+1}d$ configurazioni distinte tali che coinvolgono durante la computazione soltanto le posizioni di indice compreso tra $1-s$ e $1+s$. Ora se il costo in spazio della computazione di $\mu$ a partire da una parola di input $w$ è maggiorato da $s$ (con $s$ più grande della lunghezza $|w|$) allora tutte le configurazioni coinvolgono posizioni di indice compreso tra $1-s$ e $1+s$. Quindi se la lunghezza della computazione fosse e solo se maggiore di $(2s+1)k^{2s+1}d$ si avrebbe necessariamente una configurazione ripetuta e perciò l'insieme delle configurazioni di $\mu$ a partire da $w$ sarebbe periodico, si avrebbe non terminazione, contraddicendo il fatto che nel caso peggiore la computazione termina in un tempo finito (dalla definizione di decidibilità). Quindi il tempo di arresto è certamente maggiorato da $(2s+1)k^{2s+1}d$ che per una costante $c$ abbastanza grande indipendente da $s$ si ha che il tempo di arresto è maggiorato da $2^{cs}$.
\end{proof}
\end{prop}



    
